We introduce CPTC, a novel conformal prediction algorithm that utilizes the state prediction from SDS to improve uncertainty quantification for time series forecasting.
\vspace{-0.5em}
\paragraph{Problem Setting}
% Consider a sequential prediction problem over data generated by the SDS dynamics as illustrated in section \ref{background:sds}, whose observations $o_t \in \R^d$ for all $t$. 
We are given a sequence of observations $o_t \in \R^d$ for all $t$ potentially with  change points.
To be consistent with the standard notation in conformal prediction literature, let 
$\mathcal{X} \subseteq \R^{m}$ denote the input feature domain and $\mathcal{Y} \subseteq \R^d$ the target observation space. $ \mathcal{Z} = \{1, \dots, K\} \subset \mathbb{Z}^+$ indexes one of $K$ underlying states. The input features can include past observations and other environmental covariates.  In our implementations, $m = k \cdot d$ with $k$ being the look back window.

We consider the online conformal prediction setting where we have a sequence of data points from $\mathcal{X} \times \mathcal{Y}$: $(x_1, y_1), $ $(x_2, y_2), \ldots, (x_T , y_T )$. We have a trained state predictor for $\hat{p}(z_t|x_{0:t})$ , a state-conditioned forecaster $\hat{f}:\gX \times \gZ \to \gY$ (a probabilistic SDS model can serve both purposes), and are given a desired confidence level $1-\alpha$. At each time step $t$, we have seen $(x_1, y_1),(x_2, y_2), \ldots,(x_{t-1} , y_{t-1} )$ and the next input feature vector $x_t$. We want to produce an adaptive, time-dependent prediction interval based on input $\Gamma_t(x_t) \subseteq \mathcal{Y}$ for the unseen true target $y_t$, which we will learn at the next time step. Let $z_t \in \gZ$ be the true (latent switching) state at time $t$. %We assume that we have access to a model for probabilistic forecast $\hat{p}(z_t=z|x_{0:t})$ of the latent state at each timestep. %nonconformity score function $A(x, y): \mathcal{X} \times \mathcal{Y} \rightarrow \R_{\geq 0}$   kkfldldldfk;sdf%, denoted as $\hat{z_t}$, and a dynamics forecaster model for each mode $ \hat{f}_z$. 
Our goal is to produce prediction intervals with correct empirical coverage, i.e.    
\vspace{-0.2em}
\begin{equation}
    P (y_t \in \Gamma_t(x_t)) \geq 1-\alpha \; \text{ for any } t\geq 0
    \label{eq:goal}
\end{equation}
\vspace{-2em}
\subsection{Algorithm}
Existing online conformal prediction algorithms often rely on assumptions about how distribution would shift in time series. This is infeasible for data with change points due to  unknown shift mechanism. The main contribution our algorithm makes towards addressing this issue is \textit{anticipating} this type of abrupt shifts. 
Our CPTC algorithm is outlined in pseudo-code in algorithm \ref{alg:cptc}. In this section, we discuss the general procedure and key components of the algorithm. 


%\begin{wrapfigure}{r}{0.5\textwidth}
% \begin{wrapfigure}{L}{0.6\textwidth}
% \begin{minipage}{0.6\textwidth}

\begin{algorithm}[H]%[tb]
   \caption{Conformal Prediction for Time series with Change points (CPTC)}
   \label{alg:cptc}
   \SetAlgoLined
   \KwIn{nonconformity score function $A$, probabilistic state model $\hat{p}(z_t|x_{0:t})$ with prior $\hat{p}(z_0)$, forecaster model $\hat{f}: \mathcal{X}\times \mathcal{Z} \rightarrow \mathcal{Y} $, confidence level $1- \alpha$, learning rate $\gamma$}
   \KwOut{Prediction intervals $\Gamma(x_t)$ for $t \geq 1$}
   
   \For{$z \in \mathcal{Z}$}{
       Initialize $\gS_z \leftarrow \{\}$,  $\alpha_{z,0} \leftarrow \alpha$ 
   }
   \If{exists warm start data $\{(x_t, y_t)\}_{t=-w}^{0}$}{
       \For{$t \in [-w,0]$}{
        $\gS_{z_t} \leftarrow \gS_{z_t} \cup \{A(x_t, y_t)\}$, where $z_t \sim \hat{p}(z_t = z |x_{0:t})$; \tcp{warm start scores} 
       }
   }
   \For{$t \in [1,T]$}{
       \For{$z \in \mathcal{Z}$}{
           $\Gamma_{z,t}(x_t) \leftarrow \{ y : A(x_t, y) \leq Q^{1-\alpha_{z,t}}(\gS_z \cup \{\infty\})\}   $; \tcp{state-specific CP} %Conformal prediction interval with forecaster $\hat{f}(x_t,z)$, confidence level $1-\alpha_{z,t}$, calibrated on $\gS_z$\;
       }
       $\Gamma_t(x_t) \leftarrow$ Aggregate the $\Gamma_{z,t}(x_t)$s by \Eqref{eq:aggregation-average} or \Eqref{eq:aggregation-union}\;
       \KwOut{$\Gamma_t(x_t)$}
       Sample state $\hat{z}_t \sim \hat{p}(z_t|x_{0:t})$; \tcp{state-specific coverage target tracking}
       Update $\alpha_{\hat{z}_t, t+1} \leftarrow \alpha_{\hat{z}_t, t} + \gamma \cdot (\alpha - err_t)$, where $err_t = \mathbbm{1}\{y_t \notin \Gamma_t(x_t)\}$\;
       Update scores $\gS_{\hat{z}_t} \leftarrow \gS_{\hat{z}_t} \cup \{A(x_t, y_t)\}$\;
   }     
   \end{algorithm}
  %   \end{minipage}
  % \end{wrapfigure}



\vspace{-0.5em}
\paragraph{Switching Behavior and SDS Integration.}
We explicitly model the abrupt shift in dynamics with the switching dynamical systems (SDS) to anticipate the change. In our formulation, the latent state \(z_t\) indicates which dynamical regime (out of \(K\) possible regimes) is currently active at time \(t\). Given the SDS formulation (e.g. \Eqref{eq:sds}), we can factor the state transition probability from the conformal prediction process and do calibration for each of the $K$ dynamics separately. The goal in \Eqref{eq:goal} can therefore be written as:
%
\begin{equation}
    \sum_{z\in\gZ} P(y_t \in \Gamma_t(x_t)| z_t=z) \; P(z_t = z|x_{0:t}) \geq 1-\alpha 
\label{eq:factored_goal}
\end{equation}
\vspace{-1.2em}
\paragraph{State Prediction.} The SDS model provides us with the state predictor \(\hat{p}(z_t \mid x_t)\). In practice, SDS models \cite{ansari2021deep, linderman2016recurrent} have existing $z_t$ forecasts that we can extract. For other model architectures, one can modify them to output an extra covariate prediction for the state, or add an auxiliary model to classify or cluster the input space into different regimes. In practice, $z_t$ doesn’t necessarily have to come from a model, and can be any discrete contextual variable such as weekday/weekend/holiday or day/night.

\vspace{-0.5em}
\paragraph{Initialization and Warm Start,} For each switching state $z \in Z$, we initialize (1) a mode-specific set of nonconformity scores $\gS_z$ and (2) the mode-specific confidence level $\alpha_{z,0} = \alpha$. If there are $w$ steps of warm-start data $\{(x_t,y_t)\}_{t=-w}^0$, we initialize the sore sets $\gS_z$ by sampling $z_t \sim \hat{p}(z_t = z |x_{0:t})$ for $t= -w, \ldots, 0$, and inserting $A(x_t,y_t)$ into $\gS_{z_t}$. Depending on the size of the warm-start window and desired properties, practitioners can adjust the warm-starting approach. For example, using the entire warm-start window $\gS_{z} = \{A(x_t,y_t)\}_{t=-w}^0$  for all $z \in \gZ$  will result in better stability but less adaptability to modes. The approach does not change the algorithm's coverage guarantees. 

\vspace{-0.5em}
\paragraph{State-Specific Prediction Intervals.}  At every time step $t\in [1, T]$, CPTC creates state-specific conformal prediction intervals
for all states with nonzero probability. This allows us to generate predictions tailored to the current anticipated dynamics. Specifically, for every state $z\in \gZ$ where $\hat{p}(z|x_{0:t}) > 0$, we first obtain the point prediction through forecaster model $\hat{f}(x_t, z)$, and then the state-specific prediction interval $\Gamma_{z,t}(x_t)$ by calibrating on nonconformity scores set $\gS_z$ to the adaptive confidence level of $1-\alpha_{z,t}$.


\vspace{-0.5em}
\paragraph{Online Conformal Prediction.}
The exchangeability assumption does not apply to the online setting of our work. If we naively calibrate with online data, when the data distribution changes, the coverage probability may deviate from the target level $1-\alpha$. Therefore, we follow online CP methods such as ACI \citep{gibbs2021adaptive} to address this problem by continuously updating $\alpha_{z,t}$ to track a surrogate miscoverage rate $\alpha_{z,t}^*$ as defined in \eqref{eq:alpha_star} (for each state $z$ in our setting), an internal estimate of the target coverage.
%
\begin{equation}
    \alpha^*_{z,t} = \sup\{\beta \in [0, 1] : M_{z,t}(\beta) \leq \alpha\}, \; M_{z,t}(\alpha) = p(A(x_t, y_t) > Q^{1 - \alpha}(\gS_z))  
    \label{eq:alpha_star}
\end{equation}
%
In \eqref{eq:alpha_star} $M_{z,t}$ measures the probability that the true label $y_t$ falls outside the prediction set $\Gamma_{z,t}$ when the predicted state is $z$. %$M_t(\alpha)$ may not be equal to $\alpha$ due to distribution shifts. 
We track $\alpha^*_{z,t}$ over time by updating the estimate $\alpha_{z,t}$ with online optimization. In our implementation, we use the simple update from ACI (\Eqref{eq:aci_update}), though more sophisticated approaches such as \citep{angelopoulos2024pid} may be used in its place without affecting the overall algorithm.
\begin{equation}
\alpha_{\hat{z}_t, t+1} \leftarrow \alpha_{\hat{z}_t, t} + \gamma \cdot (\alpha - err_t), \; \text{where } err_t = \mathbbm{1} \{y_t \notin \Gamma_t(x_t) \}
\label{eq:aci_update}
\end{equation}
Here $\gamma > 0$ is the step size hyperparameter, and $\hat{z}_t \in \gZ, \; \hat{z}_t \sim \hat{p}(z_t = z |x_{0:t})$ is the state at time $t$ sampled to be updated. The update rule increases $\alpha_{\hat{z}_t, t}$ if coverage is too conservative and decreases it otherwise, ensuring that for each $z \in \gZ$, $\alpha_{z,t}$ converges to a value that maintains long-term coverage. 


\vspace{-0.5em}
\paragraph{Aggregation.} We aggregate the prediction intervals for each state $\Gamma_{z,t}(x_t)$ into one final set $\Gamma_{t}(x_t)$ using the weighted level set, as illustrated in \Eqref{eq:aggregation-average}. %The weighted approach results in sharp prediction intervals, but may risk slightly lower coverage if the classifier for $z$ is imperfect.
%
\begin{equation}
\Gamma_t(x_t) :=  \bigl\{ y:  \sum_{z \in \mathcal{Z}}
\hat{p}\bigl(z_t = z \mid x_t\bigr)
\mathbbm{1}\{\,y \in \Gamma_{z,t}(x_t)\} 
\geq 1 - \alpha \bigr\}
\label{eq:aggregation-average}
\end{equation}

\vspace{-0.8em}
Solution to the constraint in \Eqref{eq:aggregation-average}  is the minimal set that achieves the marginal coverage. It can be obtained through discretizing $\mathcal{Y}$ into a fine grid and calculating probability mass at each point, which becomes computationally expensive as the grid resolution or dimensionality increases. For faster computation and in implementation, we approximate the weighted level‐set in \Eqref{eq:aggregation-average} by taking the union of intervals of the most probable states as in \Eqref{eq:aggregation-union}. The two aggregation strategies achieve similar results on all our datasets (Appendix \ref{app:more_experiments}).
%
%Alternatively, we can opt for a more conservative aggregation by taking the union of intervals of the probable states as in  \sofi{explain why they are not far, explain approx eq}
%Alternatively, if finding the set in Eq. \ref{eq:aggregation-average} is challenging or there are high uncertainty for the state classifier, we can opt for a more conservative relaxation. Aggregation by taking the union (Eq. \ref{eq:aggregation-union}) results in larger sets and guaranteed validity over time. 
 \begin{equation}
    \Gamma_t(x_t) \approx \cup_{z\in \mathcal{Z}'_t} \Gamma_{z,t}(x_t), \; 
    \text{where } \mathcal{Z}'_t
=\;\underset{S\subseteq\mathcal Z}{\arg\min}\;\bigl\lvert S\bigr\rvert
\;\text{s.t.}\;\sum_{z\in S}p(z_t=z\mid x_t)\;\ge\;1-\alpha.
    \label{eq:aggregation-union}
\end{equation}
%
\vspace{-1.5em}
\paragraph{Modularity and the role of the state model.} Our algorithm can be viewed as running multiple instances of online conformal inference (lines 14-16 of algorithm \ref{alg:cptc}), one for each underlying state, and adaptively aggregating the confidence intervals when the underlying mode is switching. This differs from the setting where the temporal correlation between residuals needs to be learned online \citep{gibbs2024conformal, bastani2022practical, xu2023sequential, auer2023hopcpt}; the state model allows us to leverage training data for such information. Our algorithm is modular - the state model, forecaster model, and the online adaptive conformal prediction algorithms all operate independently of each other. For example, the state-specific online coverage tracking (line 15 of algorithm \ref{alg:cptc}) can be replaced with other online conformal prediction variants for data-specific applications. Our theoretical guarantees are agnostic to the forecaster models used. %If the forecasters and the state model are re-trained, our algorithm will adapt to the new model while still achieving asymptotic coverage. Section \ref{sec:theory} will elaborate on this claim. 



\subsection{Theoretical analysis}
\label{sec:theory}

In this section, we discuss theoretical guarantees of the CPTC algorithm. In particular, we establish finite-sample validity under exchangeability, asymptotic coverage in dynamic (potentially nonstationary) settings, robustness to imperfect state classification, and faster adaptation to distribution shifts that align with model-predicted state changes. See Appendix \ref{app:proofs} for proofs and theoretical details.

\vspace{-0.5em}
\paragraph{Finite-sample validity under exchangeability.}
We start by showing that when data is exchangeable (e.g. independent and identically distributed), CPTC achieves finite sample validity. The proof is standard based on showing the marginal coverage of split conformal prediction. 

\begin{proposition}[Finite-sample validity under exchangeability] If the data $(x_t, y_t)$, $t \geq 1$ are exchangeable, prediction intervals obtained via Algorithm \ref{alg:cptc} satisfy
\( P(y_t \in \Gamma_t(x_t)) \geq 1-\alpha\).
\label{prop:exchangable}
\end{proposition}

\vspace{-0.5em}
\paragraph{Asymptotic validity without assumptions.}
%\vspace{-0.3em}
Theorem \ref{thm:asymptotic_validity} ensures that CPTC provides reliable uncertainty quantification in the long run \textit{without assumptions on time-series stationarity or accurate state predictions}. (We do assume that the distribution of states and state predictions exhibit stable long-term behavior in Assumption \ref{assump:stationary_distribution}.) We achieve our desiderata of \Eqref{eq:goal} on average as $T$ grows asymptotically, consistent with the results of other online conformal prediction works. 

%
\begin{theorem}[Asymptotic validity of CPTC] \label{thm:asymptotic_validity}
For any sample size \( T \geq 1 \), the CPTC algorithm (with weighted average aggregation) achieves:
\begin{equation*}
    \left| \frac{1}{T} \sum_{t=1}^T \mathbb{E}[err_t] - \alpha \right| \leq  \frac{1}{T} \sum_{z=1}^K  \frac{ 1 - (1-c\gamma)^{|\mathcal{T}_z|}  }{\gamma} |\alpha - \alpha^{*}_z|
    \label{eq:thm1}
\end{equation*}
Where:
\vspace{-4mm}
\begin{itemize}  \setlength{\itemsep}{0pt plus 1pt}
    \item \(\text{err}_t = \mathbbm{1}[Y_t \notin \Gamma_{t}(x_t)]\), the miscoverage rate. 
    \item \( \mathcal{T}_z = \{ t \in \{1, \dots, T\} : \hat{z}_t = z \} \), the set of time steps the predicted state is \( z \).
    \item $\alpha^*_z$ is the optimal miscoverage target for timesteps $\mathcal{T}_z$.
    \item  $c$ is a miscoverage constant in Lemma \ref{lemma:aci_within_states} \cite{gibbs2021adaptive}. 
\end{itemize}
\vspace{-0.5em}
Note the RHS of \Eqref{eq:thm1} decays at a rate of $\mathcal{O}(T)$ and satisfies
 \(
 \lim_{T \to \infty} \left| \frac{1}{T} \sum_{t=1}^T \mathbb{E}[err_t] - \alpha \right|= 0
 \).
\end{theorem}

%Importantly, the algorithm adapts dynamically to stabilize its coverage as the time horizon increases.

\vspace{-0.5em}
\paragraph{Robustness to imperfect state prediction.}
In real-world applications, the state predictor  $\hat{p}(z_t|x_{0:t})$ that plays a large role in our algorithm is often imperfect. The CPTC algorithm accommodates these scenarios by tracking $\alpha^*_{z}$ for each \textit{predicted} mode $z$ with adaptive online updates. Theorem \ref{thm:imperfect_state_predictions} quantifies the effect of imperfect state prediction by showing that the miscoverage rate is bounded by the product of the misclassification rate and state-specific deviations. 
%
\vspace{0.1em}
\begin{theorem}[Finite-Sample Miscoverage Bound with Imperfect State Predictions]\label{thm:imperfect_state_predictions}
For any sample size \( T \geq 1 \), the CPTC algorithm ensures that:
\vspace{-0.2em}
\[
\left| \frac{1}{T} \sum_{t=1}^{T} \mathbb{E}[err_t]  - \alpha \right| \leq \epsilon \cdot \max_z \delta_{z,T}
\]

\vspace{-1em}
Where \( \epsilon = P(\hat{z}_t \neq z_t) \) is the error rate of the state predictions and \( \delta_{z,T} \) is the miscoverage deviation from \( \alpha \) for any predicted state $z$ defined in Lemma \ref{lemma:aci_within_states}. For all $z$, \( \delta_{z,T} \to 0 \) as \( |T| \to \infty \).
%\vspace{-4mm}
% \begin{itemize}\setlength{\itemsep}{0pt plus 1pt}
    %\item \(\text{err}_t := \mathbb{I}\big(Y_t \notin \Gamma_{CPTC}(x_t)\big)\), the miscoverage rate. 
    %\item \( \mathcal{T}_z = \{ t \in \{1, \dots, T\} : \hat{z}_t = z \} \), the set of times when the predicted state is \( z \).
    %\item \( z_t \) denotes the true latent state at time \( t \), and \( \hat{z}_t \) the predicted state. 
%     \item \( \epsilon = P(\hat{z}_t \neq z_t) \) is the error rate of the state predictions.
%     \item \( \delta_{z,T} \) is the miscoverage deviation from \( \alpha \) for any predicted state $z$ defined in Lemma \ref{lemma:aci_within_states}. For all $z$, \( \delta_{z,T} \to 0 \) as \( |T| \to \infty \).
% \end{itemize}
\end{theorem}

Theorem \ref{thm:imperfect_state_predictions} demonstrates that CPTC can maintain asymptotic valid coverage even when state predictions are imperfect, with the bound tightening as state predictions improve or the number of observations per state increases. In practice, improving the accuracy of the state prediction model directly enhances the coverage performance of CPTC. However, even with imperfect state predictions, the algorithm's adaptive calibration mechanism ensures that the overall coverage remains close to the desired level \(1 - \alpha\) as experiments show in Section \ref{sec:accuracy}.

%\subsubsection{Validity for union aggregation}

%As noted in the methodology section, if we use set union as the aggregation approach, CPTC will produce more conservative prediction intervals aiming for overcoverage. 

%By leveraging the union of prediction sets across all possible states, CPTC ensures that the algorithm maintains valid coverage, even in the presence of uncertainty during state transitions.

% \begin{proposition}[Conservative Coverage Guarantee of CPTC on Transition States]
% Note that \(\Gamma_t(x_t) = \bigcup_{z \in Z} \Gamma_{z,t}(x_t)\) is the prediction set produced by the CPTC algorithm at time \(t\).
% then the overall prediction set \(\Gamma_t(x_t)\) provides conservative coverage, meaning:
% \begin{equation}
% P(y_t \in \Gamma_t(x_t)) \geq 1 - \alpha
% \end{equation}
% for all \(t\).
% \end{proposition}

% This is because (1)  \(P(y_t \in \Gamma_{z,t}(x_t) | z_t = z) \geq 1 - \alpha_{z,t}\) for all \(t\) (i.e., the conditional coverage given state \(z\) is at least \(1 - \alpha_{z,t}\)), and (2)  \(\lim_{T_z \to \infty} (1/T_z) \sum_{t=1}^{T_z} \alpha_{z,t} \leq \alpha\) where \(T_z\) is the number of time steps where state \(z\) is predicted. 
\vspace{-0.3em}
\paragraph{Why CPTC? Faster convergence under distribution shifts.}

The strength of the CPTC algorithm is its fast adaptation in dynamic environments with distributional changes. When a distribution shift aligns with a predicted state transition, the structure of CPTC allows the confidence level $\alpha_{z,t}$ to converge to the new target error rate faster compared to purely online methods.

Consider a data stream with a distribution shift occurring at time $t_{{\text{shift}}}$. Let $\alpha^*_{j}$ denote the optimal target error rate for mode $j$ after the shift, and assume that the predicted state $\hat{z}_t$ correctly reflects the shift: $t_{\text{shift}}$ and $\alpha^*_{z} = \alpha^*_{j}$ for $t > t_{{\text{shift}}}$. For conciseness of notation in Theorem \ref{thm:adaptation}, let $\delta_{j,T}$ and $\delta_{ACI,T}$ denote the miscoverage rate \textit{starting from } $t_{\text{shift}}$ for CPTC and ACI respectively.

\vspace{0.1em}
\begin{theorem}[Miscoverage Ratio under State-Coincident Distribution Shift]
Under a state shift from $i$ to $j$ at time step $t_{\text{shift}}$, coinciding with predicted transition, the CPTC algorithm achieves faster convergence to the new target $\alpha^*_{j}$ compared to non-state-aware algorithm ACI at a ratio of:
\[
\frac{\delta_{j,T}}{\delta_{\text{ACI},T}} \leq \frac{|\alpha_{j, t_{\text{shift}}-1}  - \alpha^{*}_{j}|}{|\alpha_{ t_{\text{shift}}-1} - \alpha^{*}_{j}|}
\]
\label{thm:adaptation}
\end{theorem}
\vspace{-1em}
%
The numerator $|\alpha_{j, t_{\text{shift}}-1}  - \alpha^{*}_{j}| \rightarrow 0$ as $T \rightarrow \infty$, when in-state adaptation has converged; but the denominator $|\alpha_{ t_{\text{shift}}-1} - \alpha^{*}_{j}|$ remains greater than zero regardless of the length of the time series. This result highlights the advantage of segmenting nonconformity scores by state: the state-specific $\alpha_{z,t}$ aligns the target miscoverage updates with distribution shifts, allowing CPTC to quickly adjust its prediction intervals. Such accelerated adaptation allows for shorter miscoverage periods and is critical in real-world applications where rapid responses to changing conditions are necessary.
%Theorem \ref{thm:adaptation} shows that CPTC's bound is smaller by a factor that depends on the ratio of the distances of the original $\alpha$ and $\alpha_{i}^*$ from the new optimum $\alpha_{j}^*$, as well as the ratio of the segment lengths.

